\documentclass[12pt]{fphw}

\usepackage{fontspec}
\usepackage[greek,english]{babel}
\setmainfont{Palatino Linotype}

\usepackage{graphicx}
\usepackage{booktabs}
\usepackage{listings}
\usepackage{enumitem}
\usepackage{hyperref}
\usepackage{amsmath}
\usepackage{float}
\usepackage[framed]{ntheorem}
\usepackage{amssymb}
\usepackage{xcolor}
\usepackage{tabularx}
\usepackage{tikz}
\usepackage{titlesec}

\geometry{left=2.5cm,right=2.5cm,top=4.3cm,bottom=2.5cm,headheight=58pt,headsep=20pt}

\titleformat{\subsection}[block]{\bfseries}{\thesubsection.}{3mm}{}
\titleformat{\subsubsection}[block]{\bfseries\itshape}{\thesubsubsection.}{3mm}{}

\fancyfoot[RO]{\footnotesize\itshape Α. Ζαχαράκης, Ι. Οικονομίδης}

\graphicspath{{./images/}}

\newframedtheorem{frm-thm}{Theorem}

\hypersetup{
    colorlinks=true,
    linkcolor=blue,
    filecolor=magenta,
    urlcolor=cyan,
}
\urlstyle{same}

\setlist[itemize]{leftmargin=*}
\lstset{
    basicstyle=\ttfamily\small,
    breaklines=true,
    frame=single,
    backgroundcolor=\color{gray!5},
    keywordstyle=\color{blue},
    commentstyle=\color{green!50!black},
    stringstyle=\color{orange}
}

\newcommand{\en}[1]{\foreignlanguage{english}{#1}}

\renewcommand{\StudentLang}{Ομάδα Εργασίας}
\renewcommand{\CourseLang}{Μάθημα}
\renewcommand{\DateLang}{Εξάμηνο}

\title{Planent \\ \large Εφαρμογή Διαχείρισης Εκδηλώσεων και Ηλεκτρονικών Κρατήσεων στον Παγκόσμιο Ιστό}
\author{\\ \large \textbf{Ζαχαράκης Αντώνιος} \small (Α.Μ.: 1115202300047) \\ \large \textbf{Ιωάννης Οικονομίδης} \small (Α.Μ.: 1115201600120)}
\institute{Εθνικό και Καποδιστριακό Πανεπιστήμιο Αθηνών \\ Τμήμα Πληροφορικής \& Τηλεπικοινωνιών}
\class{Τεχνολογίες Εφαρμογών Διαδικτύου (YS14)}
\date{Εαρινό Εξάμηνο 2025–2026}

\begin{document}
\selectlanguage{greek}
\sloppy

\maketitle

\renewcommand{\contentsname}{Πίνακας Περιεχομένων}

\begingroup
\hypersetup{linkcolor=black}
\tableofcontents
\endgroup

\newpage

% =========================================================
% ΚΕΦΑΛΑΙΟ 1: ΕΙΣΑΓΩΓΗ
% =========================================================
\section{Εισαγωγή}
Η παρούσα εργασία υλοποιήθηκε στο πλαίσιο του μαθήματος «\textbf{Τεχνολογίες Εφαρμογών Διαδικτύου}» και είχε ως στόχο τη σχεδίαση και ανάπτυξη μιας ολοκληρωμένης διαδικτυακής εφαρμογής διαχείρισης εκδηλώσεων και ηλεκτρονικών κρατήσεων θέσεων, με την ονομασία \textbf{Planent}.

Η εφαρμογή ακολουθεί την κλασική αρχιτεκτονική \textbf{web browser / web server}: ο χρήστης αλληλεπιδρά με την υπηρεσία αποκλειστικά μέσω του web browser, ενώ η λογική και τα δεδομένα διαχειρίζονται από έναν web server που εκθέτει \textbf{REST API}. Πιο συγκεκριμένα, στο \en{\textbf{Backend}} χρησιμοποιήθηκε το \en{\textbf{Spring Boot}} (σε \en{Java}), ενώ το \en{\textbf{Frontend}} αναπτύχθηκε ως μια \en{Single-Page Application} με τη βιβλιοθήκη \en{\textbf{React}} σε \en{\textbf{TypeScript}}. Η επικοινωνία μεταξύ των δύο πραγματοποιείται μέσω του \en{REST API}, με τα δεδομένα να αποθηκεύονται και να διαχειρίζονται σε μια σχεσιακή βάση δεδομένων (\en{\textbf{SQL}}).

Όσον αφορά την ασφάλεια, όλη η επικοινωνία πελάτη-εξυπηρετητή κρυπτογραφείται μέσω \textbf{SSL} και πιστοποιείται με \textbf{JSON Web Tokens}, ενώ η εξουσιοδότηση γίνεται βάσει της δραστηριότητας και του ρόλου του εκάστοτε χρήστη.

Το συνολικό \en{project} αποτέλεσε μια απαιτητική αλλά και ιδιαίτερα επιμορφωτική διαδικασία, κατά την οποία δοκιμάστηκαν διάφορες στρατηγικές και μελετήθηκαν πολλοί τρόποι διεκπεραίωσης της άσκησης, αντλώντας πληροφορίες τόσο από τις διαφάνειες του μαθήματος όσο και από γενικότερες πηγές. Βασικός μας στόχος ήταν ο πειραματισμός και η εκμάθηση των τρόπων υλοποίησης μιας \en{full-stack} εφαρμογής, αλλά και η τελική δημιουργία μιας ολοκληρωμένης και εύχρηστης σελίδας, γρήγορης στις αποκρίσεις της και δίκαιης απέναντι σε όλους τους χρήστες. Έτσι, στήσαμε μια εφαρμογή με όλα τα απαραίτητα στοιχεία, όπως ασφάλεια, ταχύτητα, μαζική υποστήριξη και \en{smooth UX}, αφήνοντας εκτός μόνο ένα \en{payment method} (το οποίο μπορεί να ενσωματωθεί (\en{plugged in}) αν χρειαστεί), ώστε η εφαρμογή μας να είναι ουσιαστικά έτοιμη προς χρήση.
\newpage

% =========================================================
% ΚΕΦΑΛΑΙΟ 2: ΑΡΧΙΤΕΚΤΟΝΙΚΗ ΣΥΣΤΗΜΑΤΟΣ
% =========================================================
\section{Αρχιτεκτονική Συστήματος}
Το σύστημα είναι οργανωμένο σε τρία ανεξάρτητα \en{deployable} τμήματα: τη βάση δεδομένων (\en{\textbf{Database}}), το νωτιαίο άκρο (\en{\textbf{Backend}}) και τη διεπαφή χρήστη (\en{\textbf{Frontend}}), τα οποία εκτελούνται μέσα σε αυτόνομα και απομονωμένα \en{Docker containers}. Η βασική σχεδιαστική προσέγγιση της εφαρμογής είναι ότι ο έξω κόσμος έχει πρόσβαση αποκλειστικά και μόνο στο \en{Frontend}. Η επικοινωνία μεταξύ των υπόλοιπων επιμέρους τμημάτων δεν εκτίθεται απευθείας στο διαδίκτυο, αλλά πραγματοποιείται εσωτερικά μέσω απομονωμένων δικτύων Docker (\textbf{Docker networks}) τοπικά.

Κεντρικό ρόλο στην σχεδίαση αυτή παίζει ο \textbf{Nginx}, ο οποίος λειτουργεί ως \textbf{Web Server} και \textbf{Reverse Proxy} στην πλευρά του \en{Frontend}. Ο Nginx ακούει στις δημόσιες θύρες \textbf{443} (HTTPS) αλλά και \textbf{80} (HTTP), ώστε να ανακατευθύνει αυτόματα κάθε μη ασφαλή σύνδεση στην ασφαλή, όπου και εκτελείται το \textbf{SSL termination} μέσω ενός (\en{self-signed}) SSL πιστοποιητικού. Αυτή η σχεδίαση επιτρέπει στο \textbf{API} της εφαρμογής να ζει απευθείας μέσα στην ίδια θύρα του \en{Frontend}, με αποτέλεσμα, όταν η \en{React} εφαρμογή χρειάζεται δεδομένα από κάποιο \en{API Endpoint}, να στέλνει το αίτημα στο κατάλληλο endpoint με το πρόθεμα \en{\texttt{/api}} (\en{\texttt{https://localhost/api/}}). Ο \en{Nginx} αναχαιτίζει αυτό αίτημα, το αποκρυπτογραφεί (\en{SSL termination}) και το δρομολογεί άμεσα στο εσωτερικό, απομονωμένο \en{container} του \en{Backend} μέσω του απλού, μη κρυπτογραφημένου πρωτοκόλλου \en{HTTP}.

Με αυτή την αρχιτεκτονική σχεδίαση επιτυγχάνεται διπλό όφελος για την εφαρμογή. Από τη μία πλευρά, εγγυάται η μέγιστη δυνατή ασφάλεια, καθώς το \en{Backend} και η βάση δεδομένων δεν εκθέτουν καμία θύρα πρόσβασης στο εξωτερικό περιβάλλον, παραμένοντας αόρατα στον έξω κόσμο. Από την άλλη πλευρά, επειδή η επεξεργασία των πιστοποιητικών (SSL) σταματά στον \en{proxy}, το \en{Backend} απαλλάσσεται πλήρως από το υπολογιστικό βάρος της διαχείρισης των κρυπτογραφικών υπολογισμών, διατηρώντας υψηλή απόδοση και ταχύτητα απόκρισης.


\subsection{Συνοπτική Επισκόπηση Frontend}
Το \en{\textbf{Frontend}} είναι μια \en{\textbf{Single-Page Application (SPA)}} υλοποιημένη με τη βιβλιοθήκη \en{\textbf{React}} και γραμμένη σε \en{\textbf{TypeScript}}. Κατά τη φάση της παραγωγής, η εφαρμογή μεταγλωττίζεται σε βελτιστοποιημένα στατικά αρχεία (\en{HTML, CSS, JS}) μέσω του εργαλέιου \en{\textbf{Vite}} και σερβίρεται αποδοτικά από τον \en{\textbf{Nginx}}.

Η πλοήγηση πραγματοποιείται στην πλευρά του πελάτη μέσω \en{\textbf{React Router}}, αποφεύγοντας τις επαναφορτώσεις σελίδων. Η διαχείριση της κατάστασης (όπως τα στοιχεία του χρήστη, το \en{JSON Web Token} και το \en{light/dark theme}) επιτυγχάνεται μέσω του \en{\textbf{React Context API}} (\en{AuthContext} και \en{ThemeContext}). Η επικοινωνία με το REST API (\en{Backend}) υλοποιείται μέσω της βιβλιοθήκης \en{\textbf{Axios}}, που προσαρτά αυτόματα το \en{JSON Web Token} ως \en{\textbf{Bearer} token} στο \en{\textbf{Authorization Header}} κάθε εξερχόμενου αιτήματος.

\subsection{Συνοπτική Επισκόπηση Backend}
Το \en{\textbf{Backend}} της εφαρμογής έχει αναπτυχθεί με τη χρήση του \en{\textbf{Spring Boot 4}} σε \en{\textbf{Java 21}}, υιοθετώντας τη στρωματοποιημένη αρχιτεκτονική (\en{\textbf{Layered Architecture}}) για τον σαφή διαχωρισμό των αρμοδιοτήτων (\en{Separation of Concerns}). Η αρχιτεκτονική αυτή δομείται στα εξής επίπεδα:
\begin{itemize}[label=\textbullet]
    \item \en{\textbf{Controllers:}} Υπεύθυνοι για την έκθεση του \en{\textbf{REST API}} της εφαρμογής. Δέχονται τα \en{HTTP} αιτήματα (\en{GET, POST} κ.λπ.), ελέγχουν αν τα δεδομένα που στέλνει ο χρήστης είναι σωστά μέσω \en{validation} στα \en{\textbf{DTOs}} και στη συνέχεια αναθέτουν την εκτέλεση της λογικής στα \en{Services}, επιστρέφοντας την τελική απάντηση.
    \item \en{\textbf{Services:}} Ενσωματώνουν την επιχειρησιακή λογική και τους κανόνες του συστήματος (π.χ. έλεγχος διαθεσιμότητας θέσεων πριν από μια κράτηση). Επίσης, διαχειρίζονται τις συναλλαγές (\en{\textbf{transactions}}) στη βάση δεδομένων μέσω των \en{Repositories}, εξασφαλίζοντας ότι μια ενέργεια είτε θα πετύχει ολοκληρωτικά είτε θα ακυρωθεί αν προκύψει κάποιο σφάλμα (\en{rollback}).
    \item \en{\textbf{Repositories:}} Διεπαφές (\en{interfaces}) του \en{\textbf{Spring Data JPA}} οι οποίες, σε συνδυασμό με το \en{\textbf{Hibernate}} (\en{\textbf{ORM}}), επιτρέπουν την αλληλεπίδραση με τη \en{\textbf{MySQL}} βάση δεδομένων χρησιμοποιώντας αντικείμενα της \en{Java} αντί για χειροκίνητα \en{queries}. Παρέχουν έτοιμες τις βασικές λειτουργίες αποθήκευσης και διαγραφής (\en{\textbf{CRUD}}), καθώς και τη δυνατότητα για δυναμικές αναζητήσεις (dynamic queries).
    \item \en{\textbf{Models (Entities):}} Οι κλάσεις της \en{Java} που αντιστοιχούν ακριβώς στους πίνακες της βάσης δεδομένων (όπως \en{User, Event, Booking}). Ορίζουν τη δομή των δεδομένων, τις σχέσεις μεταξύ τους (\en{one-to-many, many-to-many}) και τους περιορισμούς των πεδίων (όπως είναι και στην βάση).
    \item \en{\textbf{DTOs (Data Transfer Objects):}} Ξεχωριστά μοντέλα δεδομένων που χρησιμοποιούνται αποκλειστικά για τη μεταφορά πληροφοριών στα \en{Requests} και \en{Responses}. Χρησιμοποιούνται για να μην εκτίθενται τα πραγματικά \en{entities} της βάσης στο διαδίκτυο, επιτρέποντάς μας να φιλτράρουμε την πληροφορία, στέλνοντας μόνο τα απαραίτητα πεδία στο \en{Frontend}.
\end{itemize}

Η ασφάλεια της εφαρμογής και ο έλεγχος πρόσβασης των χρηστών βασίζονται καθολικά στο \en{\textbf{Spring Security}}, το οποίο αναλαμβάνει την αρχική αυθεντικοποίηση μέσω login, την μετέπειτα \en{stateless} αυθεντικοποίηση μέσω \en{JWT} και, τελευταία, την εξουσιοδότηση βάσει ρόλων.

\begin{table}[h!]
    \centering
    \small
    \begin{tabularx}{\textwidth}{lX}
    \toprule
    \textbf{Επίπεδο} & \textbf{Τεχνολογίες} \\
    \midrule
    \en{\textbf{Frontend}} & \en{React 19, TypeScript, Vite, React Router, Axios, React-Leaflet (OpenStreetMap)} \\
    \addlinespace[4pt]
    \en{\textbf{Backend}} & \en{Java 21, Spring Boot 4, Spring Web MVC, Spring Security, Spring Data JPA, Jackson (JSON \& XML), Thumbnailator} \\
    \addlinespace[4pt]
    \en{\textbf{Database}} & \en{MySQL 9.6.0} (σχεσιακή ΒΔ) \\
    \addlinespace[4pt]
    \en{\textbf{Ασφάλεια}} & \en{SSL/TLS, JWT (jjwt), BCrypt Password Hashing, CORS, Stateless Sessions} \\
    \addlinespace[4pt]
    \en{\textbf{Build}} & \en{Maven Wrapper (Backend), npm + Vite (Frontend), Dockerfiles (Production)} \\
    \bottomrule
    \end{tabularx}
    \caption{Tech Stack της Εφαρμογής}
\end{table}

\newpage

% =========================================================
% ΚΕΦΑΛΑΙΟ 3: ΛΕΙΤΟΥΡΓΙΕΣ & ΡΟΛΟΙ ΧΡΗΣΤΩΝ
% =========================================================
\section{Λειτουργίες \& Ρόλοι Χρηστών}
Αν και στην πράξη οι χρήστες διαχωρίζονται ανάλογα με τη δραστηριότητά τους μέσα στην πλατφόρμα, σε επίπεδο κώδικα και ασφάλειας (\en{Spring Security}) το σύστημα υποστηρίζει ουσιαστικά δύο βασικούς ρόλους: τον Διαχειριστή (\en{\texttt{ADMIN}}) και τον εγγεγραμμένο Χρήστη (\en{\texttt{USER}}), πέρα φυσικά από τον μη συνδεδεμένο Επισκέπτη.

Η σχεδιαστική λεπτομέρεια της εφαρμογής είναι ότι οι ιδιότητες του «Συμμετέχοντα» και του «Διοργανωτή» δεν αποτελούν ξεχωριστούς στατικούς ρόλους στη βάση δεδομένων, αλλά, κάθε απλός εγγεγραμμένος χρήστης (\en{\texttt{USER}}) μπορεί να λειτουργήσει και με τις δύο ιδιότητες δυναμικά. Το δικαίωμα να επεξεργαστεί, να προβάλει ή να ακυρώσει μια εκδήλωση ή μια κράτηση ελέγχεται από το \en{Backend} με βάση το αν ο τρέχων χρήστης είναι ο δημιουργός της εκάστοτε οντότητας μέσω ελέγχου ιδιοκτησίας.

\subsection{Επισκέπτης (\en{Guest})}
Είναι ο μη συνδεδεμένος χρήστης της εφαρμογής. Οι δυνατότητές του περιορίζονται αποκλειστικά στην περιήγηση:
\begin{itemize}[leftmargin=*]
    \item Πρόσβαση στη σελίδα καλωσορίσματος και προβολή των προτεινόμενων εκδηλώσεων.
    \item Πλοήγηση και αναζήτηση όλων των δημοσιευμένων εκδηλώσεων με τη χρήση των προηγμένων φίλτρων.
    \item Προβολή των πλήρων λεπτομερειών μιας (\en{non-draft}) εκδήλωσης και του διαδραστικού χάρτη τοποθεσίας.
    \item \textbf{Περιορισμοί:} Δεν έχει δικαίωμα να δημιουργήσει εκδηλώσεις, να πραγματοποιήσει κρατήσεις θέσεων ή να ανταλλάξει μηνύματα.
\end{itemize}

\subsection{Εγγεγραμμένος Χρήστης (\en{User})}
Κάθε αυθεντικοποιημένος χρήστης με τον ρόλο \en{\texttt{USER}} έχει πρόσβαση στο σύνολο των λειτουργιών, οι οποίες διαμορφώνονται ανάλογα με τη δραστηριότητά του:
\begin{itemize}[leftmargin=*]
    \item \textbf{Ως Συμμετέχων (\en{Attendee}):} Από τη στιγμή που ο χρήστης πραγματοποιεί μια κράτηση εισιτηρίων για μια εκδήλωση, το \en{Backend} τον καταχωρεί ως ιδιοκτήτη της συγκεκριμένης κράτησης. Έτσι, αποκτά το δικαίωμα να τη διαχειρίζεται, να την επισκοπεί (\en{\texttt{/my-bookings}}), να την ακυρώσει (εφόσον η εκδήλωση είναι ενεργή) και να συνομιλήσει μέσω μηνυμάτων με τον διοργανωτή της εκδήλωσης.
    \item \textbf{Ως Διοργανωτής (\en{Organizer}):} Οποιοσδήποτε χρήστης μπορεί να δημιουργήσει μια νέα εκδήλωση. Μόλις η εκδήλωση αποθηκευτεί, ο χρήστης αναγνωρίζεται ως ο επίσημος διοργανωτής της και αποκτά το δικαίωμα να την επεξεργαστεί, να τη δημοσιεύσει, να την ακυρώσει, να βλέπει αναλυτικά τις κρατήσεις που κάνουν οι άλλοι χρήστες ή να στείλει μαζικό ενημερωτικό μήνυμα (\en{broadcast}) σε όλους τους συμμετέχοντες. Διαγραφή της εκδήλωσης επιτρέπεται μόνο αν βρίσκεται σε κατάσταση \en{\texttt{DRAFT}} ή δεν έχει γίνει ακόμα καμία κράτηση.
\end{itemize}

\subsection{Διαχειριστής (\en{Administrator})}
Διαθέτει τον ρόλο \en{\texttt{ADMIN}} και έχει καθολική εποπτεία πάνω σε όλη την πλατφόρμα, εκτός των μηνυμάτων, παρακάμπτοντας τους περιορισμούς ιδιοκτησίας των απλών χρηστών:
\begin{itemize}[leftmargin=*]
    \item Πρόσβαση στην αποκλειστική σελίδα διαχείρισης (\en{\texttt{/admin}}) με ξεχωριστούς πίνακες ελέγχου για χρήστες (\en{Users}) και εκδηλώσεις (\en{Events}).
    \item Έλεγχος, έγκριση ή οριστική απόρριψη (διαγραφή από τη βάση) των εκκρεμών αιτήσεων εγγραφής των χρηστών.
    \item Δυνατότητα άμεσης παρέμβασης, τροποποίησης, δημοσίευσης ή ακύρωσης οποιασδήποτε εκδήλωσης ή κράτησης στο σύστημα.
    \item Δυνατότητα εξαγωγής ολόκληρης της βάσης δεδομένων των εκδηλώσεων σε δομημένη μορφή \en{XML} ή \en{JSON}.
\end{itemize}

\newpage

% =========================================================
% ΚΕΦΑΛΑΙΟ 4: ΣΧΕΔΙΑΣΗ ΒΑΣΗΣ ΔΕΔΟΜΕΝΩΝ
% =========================================================
\section{Σχεδίαση Βάσης Δεδομένων}
Η σχεδίαση της σχεσιακής βάσης δεδομένων έγινε με γνώμονα τις λειτουργικές απαιτήσεις της εφαρμογής, αλλά και την συμβατότητα με το ζητούμενο \en{\textbf{Document Type Definition (DTD)}} για την εξαγωγή των δεδομένων. Ως Σύστημα Διαχείρισης Βάσεων Δεδομένων (\en{RDBMS}) χρησιμοποιείται η \en{\textbf{MySQL}} με την \textbf{InnoDB} ως \en{storage engine}, η οποία εξασφαλίζει την ακεραιότητα των δεδομένων μέσω ξένων κλειδιών (\en{foreign keys}), μοναδικών περιορισμών (\en{unique keys}) και βασικών κανόνων ελέγχου (π.χ. περιορισμός \en{\texttt{available <= quantity}} στα εισιτήρια). Οι πίνακες της βάσης είναι:

\vspace{-0.25cm}
\begin{table}[H]
    \centering
    \small
    \begin{tabularx}{\textwidth}{lX}
    \toprule
    \textbf{Πίνακας} & \textbf{Συνοπτική περιγραφή} \\
    \midrule
    \en{\texttt{User}} & Στοιχεία χρήστη (\en{username, password hash}, στοιχεία επικοινωνίας, διεύθυνση, γεωγραφική θέση, ΑΦΜ, ρόλος \en{admin}, κατάσταση έγκρισης). \\
    \addlinespace[1pt]
    \en{\texttt{Event}} & Εκδηλώσεις (τίτλος, τύπος, χώρος, διεύθυνση/γεωγραφικές συντεταγμένες, χρόνοι έναρξης/λήξης, χωρητικότητα, περιγραφή, κατάσταση: \en{DRAFT / PUBLISHED / COMPLETED / CANCELLED, organizer}). \\
    \addlinespace[1pt]
    \en{\texttt{Category}} & Προκαθορισμένη λίστα κατηγοριών (\en{Music, Sports, Conference}, κ.λπ.). \\
    \addlinespace[1pt]
    \en{\texttt{Event\_Category}} & Πίνακας ένωσης \en{Event} $\leftrightarrow$ \en{Category} (πολλαπλές κατηγορίες ανά εκδήλωση). \\
    \addlinespace[1pt]
    \en{\texttt{Event\_TicketType}} & Τύποι εισιτηρίων για κάθε εκδήλωση (όνομα, τιμή, ποσότητα, διαθέσιμα). \\
    \addlinespace[1pt]
    \en{\texttt{Event\_Media}} & Προαιρετικές φωτογραφίες (\en{file path} στον δίσκο) ανά εκδήλωση. \\
    \addlinespace[1pt]
    \en{\texttt{Booking}} & Κρατήσεις (\en{attendee, ticket\_type}, πλήθος εισιτηρίων, συνολικό κόστος, χρονοσήμανση, κατάσταση: \en{PENDING / CONFIRMED / CANCELLED}). \\
    \addlinespace[1pt]
    \en{\texttt{Message}} & Μηνύματα μεταξύ χρηστών (\en{sender, receiver, event}, σώμα, \en{isRead}, σημαίες λογικής διαγραφής \en{deletedBySender/deletedByReceiver}). \\
    \addlinespace[1pt]
    \en{\texttt{User\_Event\_Interaction}} & Το σκορ (\en{rating}) που προκύπτει από τις προβολές ή τις κρατήσεις σε εκδηλώσεις για κάθε χρήστη. Πίνακας εισόδου για την εκπαίδευση του αλγορίθμου συστάσεων. \\
    \addlinespace[1pt]
    \en{\texttt{User\_Recommendation\_Vector}} & Οι λανθάνοντες παράγοντες (\en{latent factors}) και το \en{bias} για κάθε χρήστη για τον υπολογισμό των προσωποποιημένων προτάσεων. \\
    \addlinespace[1pt]
    \en{\texttt{Event\_Recommendation\_Vector}} & Οι λανθάνοντες παράγοντες (\en{latent factors}) και το \en{bias} για κάθε εκδήλωση για τον υπολογισμό των προσωποποιημένων προτάσεων. \\
    \addlinespace[1pt]
    \en{\texttt{Recommendation\_Config}} & Πίνακας μιας σειράς που κρατάει το καθολικό σφάλμα (\en{global bias}) και τις υπερπαραμέτρους εκπαίδευσης (\en{learning rate, epochs, regularization}) του αλγορίθμου. \\
    \bottomrule
    \end{tabularx}
\end{table}


\subsection{Βασικοί Περιορισμοί}
\begin{itemize}
    \item \textbf{Σχέσεις Πινάκων:} Η διασύνδεση των δεδομένων βασίζεται σε περιορισμούς ξένων κλειδιών (\en{Foreign Keys}). Οι πολυμερείς σχέσεις, όπως αυτή μεταξύ εκδηλώσεων και κατηγοριών, επιλύονται με τη χρήση του πίνακα ένωσης \en{\texttt{Event\_Category}} (\en{Many-to-Many}).
    \item \textbf{\en{Cascade Rules}:} Στη διαγραφή μιας εκδήλωσης έχει οριστεί κανόνας \en{\texttt{CascadeType.ALL}} σε συνδυασμό με \en{\texttt{orphanRemoval}} σε επίπεδο \en{JPA}. Με αυτόν τον τρόπο, η διαγραφή μιας εγγραφής στον πίνακα \en{\texttt{Event}} πυροδοτεί την αυτόματη εκκαθάριση των αντίστοιχων εξαρτημένων εγγραφών στους πίνακες των πολυμέσων (\en{\texttt{Event\_Media}}), των τύπων εισιτηρίων (\en{\texttt{Event\_TicketType}}) και των κατηγοριών, αποτρέποντας την ύπαρξη ορφανών δεδομένων.
    \item \textbf{Περιορισμοί Τιμών:} Ο πίνακας των εισιτηρίων ενσωματώνει κανόνα ελέγχου ο οποίος διασφαλίζει ότι ο αριθμός των διαθέσιμων εισιτηρίων δεν μπορεί ποτέ να ξεπεράσει την αρχική ποσότητα (\en{\texttt{available <= quantity}}), λειτουργώντας ως μια επιπλέον δικλείδα ασφαλείας στο επίπεδο της βάσης.
\end{itemize}


\subsection{Αρχικοποίηση}
Κατά την πρώτη εκκίνηση του \en{Backend} (\en{\texttt{PlanentApplication\#initDatabase}}) δημιουργείται αυτόματα ένας ενσωματωμένος διαχειριστής με \en{\textbf{admin}} όνομα χρήστη και κωδικό πρόσβασης. Επιπλέον, αρχικοποιούνται αυτόματα και οι βασικές κατηγορίες εκδηλώσεων της πλατφόρμας.

\newpage

% =========================================================
% ΚΕΦΑΛΑΙΟ 5: BACKEND — REST API
% =========================================================
\section{Backend — REST API}
Το \en{\textbf{Backend}} εκθέτει ένα \en{\textbf{RESTful API}} με αρχικό \en{context path} \en{\textbf{/api}} που καταναλώνεται από το \en{Frontend}. Όλα τα αιτήματα προωθούνται εσωτερικά από το Frontend μέσω του \en{Nginx Reverse Proxy} και, εκτός από τα δημόσια endpoints, απαιτούν την ύπαρξη ενός έγκυρου \en{\textbf{JWT}} στο \en{Authorization Header} ως \en{Bearer token}. Το \en{\textbf{JSON Web Token}} κουβαλάει τις ακόλουθες πληροφορίες (\en{claims}) για τον χρήστη:
\begin{itemize}
    \item \en{\textbf{User ID} (\texttt{userId}):} Το μοναδικό αναγνωριστικό (\en{ID}) του χρήστη στη βάση δεδομένων.
    \item \en{\textbf{Username} (\texttt{subject}):} Το μοναδικό όνομα χρήστη (\en{username}) με το οποίο έγινε η σύνδεση.
    \item \en{\textbf{Role} (\texttt{isAdmin}):} Μια \en{boolean} τιμή που δείχνει αν ο χρήστης είναι διαχειριστής ή όχι, η οποία μεταφράζεται σε δικαιώματα \en{\texttt{ADMIN}} ή \en{\texttt{USER}} (\en{Authorities}) μέσα στην εφαρμογή. Αυτό επιτρέπει στο σύστημα να αναγνωρίζει τον βασικό ρόλο χωρίς επιπλέον κλήση στη βάση δεδομένων.
\end{itemize}

Η διαδικασία αυθεντικοποίησης ελέγχεται αυτόματα σε κάθε αίτηση μέσω του φίλτρου \en{\textbf{JwtAuthFilter}}, το οποίο επικυρώνει την υπογραφή του \en{token} και φορτώνει τα παραπάνω βασικά στοιχεία του χρήστη στο \en{Security Context} του \en{Spring Boot}.

\subsection{Συνοπτικοί Πίνακες Endpoints}

\subsubsection{Αυθεντικοποίηση \& Διαχείριση Χρηστών (/auth , /users)}
\begin{table}[h!]
    \centering
    \small
    \begin{tabularx}{\textwidth}{llX}
    \toprule
    \textbf{Method} & \textbf{Endpoint} & \textbf{Περιγραφή} \\
    \midrule
    \en{POST} & \en{\texttt{/auth/register}} & Εγγραφή νέου χρήστη στο σύστημα (\en{public}, απαιτεί έγκριση από \en{admin}). \\
    \addlinespace[1pt]
    \en{POST} & \en{\texttt{/auth/login}} & Σύνδεση χρήστη και έκδοση του \en{JWT} (\en{public}). \\
    \addlinespace[1pt]
    \en{GET} & \en{\texttt{/users/me}} & Επιστροφή των στοιχείων προφίλ του τρέχοντος χρήστη. \\
    \addlinespace[1pt]
    \en{GET} & \en{\texttt{/users/\{userId\}}} & Προβολή προφίλ συγκεκριμένου χρήστη (\en{owner} ή \en{admin}). \\
    \addlinespace[1pt]
    \en{PATCH} & \en{\texttt{/users/\{userId\}}} & Ενημέρωση στοιχείων χρήστη (\en{owner} ή \en{admin}). \\
    \addlinespace[1pt]
    \en{GET} & \en{\texttt{/users}} & Λίστα όλων των εγγεγραμμένων χρηστών (\en{admin only}). \\
    \addlinespace[1pt]
    \en{PATCH} & \en{\texttt{/users/\{userId\}/approve}} & Έγκριση αίτησης εγγραφής διοργανωτή (\en{admin only}). \\
    \addlinespace[1pt]
    \en{DELETE} & \en{\texttt{/users/\{userId\}}} & Απόρριψη αίτησης ή διαγραφή χρήστη (\en{admin only}). \\
    \bottomrule
    \end{tabularx}
\end{table}

\subsubsection{Διαχείριση Εκδηλώσεων (/events)}
\begin{table}[H]
    \centering
    \small
    \begin{tabularx}{\textwidth}{llX}
    \toprule
    \textbf{Method} & \textbf{Endpoint} & \textbf{Περιγραφή} \\
    \midrule
    \en{GET} & \en{\texttt{/events}} & Σελίδες προτεινόμενων και ορατών εκδηλώσεων (\en{public}). \\
    \addlinespace[1pt]
    \en{GET} & \en{\texttt{/events/\{eventId\}}} & Προβολή εκδήλωσης (\en{public}, εμφανίζει \en{drafts} μόνο στον \en{organizer}). \\
    \addlinespace[1pt]
    \en{GET} & \en{\texttt{/events/search}} & Σελιδοποιημένη αναζήτηση προτεινώμενων εκδηλώσεων με φίλτρα (\en{public}). \\
    \addlinespace[1pt]
    \en{GET} & \en{\texttt{/events/categories}} & Επιστροφή όλων των διαθέσιμων κατηγοριών εκδηλώσεων (\en{public}). \\
    \addlinespace[1pt]
    \en{GET} & \en{\texttt{/events/my-events}} & Σελίδες των εκδηλώσεων του συνδεδεμένου διοργανωτή. \\
    \addlinespace[1pt]
    \en{POST} & \en{\texttt{/events}} & Δημιουργία εκδήλωσης με κείμενο και αρχεία πολυμέσων. \\
    \addlinespace[1pt]
    \en{PATCH} & \en{\texttt{/events/\{eventId\}}} & Ενημέρωση στοιχείων και εικόνων εκδήλωσης (\en{organizer} ή \en{admin}). \\
    \addlinespace[1pt]
    \en{DELETE} & \en{\texttt{/events/\{eventId\}}} & Διαγραφή εκδήλωσης από τη βάση δεδομένων (\en{organizer} ή \en{admin}). \\
    \addlinespace[1pt]
    \en{POST} & \en{\texttt{/events/\{eventId\}/view}} & Καταγραφή προβολής εκδήλωσης για το σύστημα συστάσεων. \\
    \bottomrule
    \end{tabularx}
\end{table}

\subsubsection{Διαχείριση Κρατήσεων (/bookings)}
\begin{table}[h!]
    \centering
    \small
    \begin{tabularx}{\textwidth}{llX}
    \toprule
    \textbf{Method} & \textbf{Endpoint} & \textbf{Περιγραφή} \\
    \midrule
    \en{GET} & \en{\texttt{/bookings/my-bookings}} & Προβολή των κρατήσεων του συνδεδεμένου χρήστη. \\
    \addlinespace[1pt]
    \en{GET} & \en{\texttt{/bookings/event/\{eventId\}}} & Σελίδες κρατήσεων μιας εκδήλωσης (\en{organizer} ή \en{admin}). \\
    \addlinespace[1pt]
    \en{GET} & \en{\texttt{/bookings/\{bookingId\}}} & Ανάκτηση συγκεκριμένης κράτησης βάσει \en{ID} (\en{attendee} ή \en{admin}). \\
    \addlinespace[1pt]
    \en{POST} & \en{\texttt{/bookings}} & Δημιουργία νέας κράτησης εισιτηρίων με \en{atomic} έλεγχο. \\
    \addlinespace[1pt]
    \en{POST} & \en{\texttt{/bookings/\{bookingId\}/cancel}} & Ακύρωση υπάρχουσας κράτησης (\en{attendee} ή \en{admin}). \\
    \bottomrule
    \end{tabularx}
\end{table}

\newpage

\subsubsection{Σύστημα Μηνυμάτων \& Εξαγωγές (/messages , /admin/export)}
\begin{table}[h!]
    \centering
    \small
    \begin{tabularx}{\textwidth}{llX}
    \toprule
    \textbf{Method} & \textbf{Endpoint} & \textbf{Περιγραφή} \\
    \midrule
    \en{GET} & \en{\texttt{/messages/inbox}} & Ανάκτηση των εισερχόμενων μηνυμάτων του χρήστη. \\
    \addlinespace[1pt]
    \en{GET} & \en{\texttt{/messages/sent}} & Ανάκτηση των απεσταλμένων μηνυμάτων του χρήστη. \\
    \addlinespace[1pt]
    \en{GET} & \en{\texttt{/messages/unread-count}} & Επιστροφή του αριθμού των μη αναγνωσμένων μηνυμάτων. \\
    \addlinespace[1pt]
    \en{POST} & \en{\texttt{/messages/broadcast/\{eventId\}}} & Μαζική αποστολή μηνύματος σε συμμετέχοντες (\en{organizer} ή \en{admin}). \\
    \addlinespace[1pt]
    \en{GET} & \en{\texttt{/messages/\{messageId\}}} & Προβολή συγκεκριμένου μηνύματος βάσει \en{ID}. \\
    \addlinespace[1pt]
    \en{PATCH} & \en{\texttt{/messages/\{messageId\}/read}} & Σήμανση ενός μηνύματος ως αναγνωσμένο (\en{receiver only}). \\
    \addlinespace[1pt]
    \en{POST} & \en{\texttt{/messages}} & Αποστολή νέου προσωπικού μηνύματος στον διοργανωτή μιας εκδήλωσης (απαιτεί κράτηση στην εκδήλωση). \\
    \addlinespace[1pt]
    \en{DELETE} & \en{\texttt{/messages/\{messageId\}}} & Λογική διαγραφή μηνύματος από τον χρήστη. \\
    \midrule
    \en{GET} & \en{\texttt{/admin/export/events}} & Εξαγωγή όλων των εκδηλώσεων σε \en{XML} ή \en{JSON} (\en{admin only}). \\
    \bottomrule
    \end{tabularx}
\end{table}


\subsection{Υλοποίηση Επιχειρησιακής Λογικής (\en{Services})}
Η λογική της εφαρμογής βρίσκεται στο επίπεδο των υπηρεσιών (\en{Services}), όπου τηρούνται όλοι οι επιχειρησιακοί κανόνες και οι μηχανισμοί διασφάλισης των δεδομένων:

\subsubsection{Διαχείριση Εκδηλώσεων (\en{Events})}
\begin{itemize}[leftmargin=*]
    \item \textbf{Κύκλος ζωής:} Για την αυτόματη ενημέρωση των εκδηλώσεων, χρησιμοποιείται ένας background scheduler (\en{@Scheduled}) που εκτελείται κάθε 15 λεπτά. Αυτός τρέχει ένα query στη βάση που βρίσκει όσα \en{PUBLISHED} events έχουν λήξει με βάση την τρέχουσα ώρα και αλλάζει την κατάστασή τους σε \en{COMPLETED}.
    \item \textbf{Περιορισμοί κατάστασης:} Δεν επιτρέπεται να γίνει τροποποίηση σε εκδήλωση που είναι ήδη ακυρωμένη ή ολοκληρωμένη, ούτε και επιτρέπεται να διαγραφτεί ένα δημοσιευμένο event αν αυτό έχει ήδη συγκεντρώσει έστω και μία κράτηση.
    \item \textbf{Επικύρωση κατά create/update:} Κατά τη διαχείριση μιας εκδήλωσης, διασφαλίζεται πως όλα τα δεδομένα είναι έγκυρα πριν αποθηκευτούν στην βάση δεδομένων. Συγκεκριμένα ελέγχεται ότι το άθροισμα των ποσοτήτων όλων των επιμέρους τύπων δεν ξεπερνά τη συνολική χωρητικότητα του χώρου, οτι ο διοργανωτής δεν μειώνει την ποσότητα κάτω από τον αριθμό των θέσεων που έχουν ήδη πουληθεί κ.α.. Αυτο επιτυγχάνεται στο τέλος του create και του update, όπου καλείται η κεντρική μέθοδος \en{event.validate()}, η οποία ελέγχει συνολικά την οντότητα (ημερομηνίες, χωρητικότητες, ύπαρξη εισιτηρίων και κατηγοριών για δημοσιευμένα events) ώστε να διασφαλιστεί ότι τα πάντα έχουν ενημερωθεί σωστά και είναι έγκυρα πριν το τελικό save. Να σημειωθεί πως τα draft events μπορουν να εχουν ελλιπη πεδία κατά την ενημέρωση τους, ωστόσο, πρέπει να συμπληρωθούν για να επιτραπεί η δημοσίευση τους.
    \item \textbf{Σύστημα αρχείων και rollback:} Κατά την επεξεργασία των πολυμέσων, τα αρχεία αποθηκεύονται στον δίσκο, ο οποίος όμως δεν υποστηρίζει transactional rollback όπως η βάση δεδομένων. Για τον λόγο αυτό, οι διαδρομές των αρχείων που ανεβαίνουν καταγράφονται σε μια λίστα και, αν υπάρξει οποιοδήποτε σφάλμα πριν το τελικό save στη βάση, ένα \en{catch} block σβήνει αμέσως τα αρχεία από το filesystem για να μην μείνουν ορφανά.
    \item \textbf{Υλοποίηση συστάσεων:} Οι υπολογισμοί του αλγορίθμου \en{Matrix Factorization} έχουν ενσωματωθεί απευθείας στα SQL queries της βάσης. Για τη διαχείριση του προβλήματος του \en{cold start} ή των επισκεπτών, χρησιμοποιείται η SQL συνάρτηση \en{COALESCE}, η οποία μηδενίζει τους λανθάνοντες παράγοντες και επιστρέφει ως baseline σκορ μόνο το γενικό bias (τη δημοτικότητα) της εκδήλωσης.
    \item \textbf{Caching:} Επειδή όλοι οι επισκέπτες βλέπουν τις ίδιες προτεινόμενες εκδηλώσεις, αυτές αποθηκεύονται στην τοπική προσωρινή μνήμη (cache) για την άμεση πρόσβαση τους, χωρίς queries στην βάση. Αυτή η cache αδειάζει καθολικά κάθε φορά που εκτελείται δημιουργία, τροποποίηση, διαγραφή ή η αυτοματοποιημένη ολοκλήρωση μιας εκδήλωσης προκειμένου να διατηρηθεί η εγκυρότητα/επικαιροποίηση των δεδομένων.
\end{itemize}

\subsubsection{Διαχείριση Κρατήσεων (\en{Bookings})}
\begin{itemize}[leftmargin=*]
    \item \textbf{Δημιουργία κρατήσεων:} Κατά την αγορά εισιτηρίων, προκειμένου να αποφευχθούν συνθήκες ανταγωνισμού (\en{race conditions}) και υπερκρατήσεις, το σύστημα εφαρμόζει απαισιόδοξο κλείδωμα (\en{pessimistic lock}) στην εγγραφή του συγκεκριμένου τύπου εισιτηρίου στη βάση. Ακόμη, απαγορεύεται στους διοργανωτές να κλείνουν εισιτήρια για τις δικές τους εκδηλώσεις.
    \item \textbf{Ακύρωση κρατήσεων:} Παρόμοια και εδω κλειδώνεται η εγγραφή του εισιτηρίου πρώτα πριν ακυρωθεί η κράτηση και έπειτα επιστρέφονται οι θέσεις στο διαθέσιμο απόθεμα. Δεν επιτρέπεται ακύρωση σε COMPLETED events.
\end{itemize}

\subsubsection{Σύστημα Μηνυμάτων (\en{Messages})}
\begin{itemize}[leftmargin=*]
    \item \textbf{Κανόνες επικοινωνίας:} Επιτρέπεται η ανταλλαγή μηνυμάτων αποκλειστικά ανάμεσα στον διοργανωτή μιας εκδήλωσης και σε έναν χρήστη που έχει ενεργή κράτηση σε αυτήν, ενώ απαγορεύεται η αποστολή μηνύματος στον εαυτό μας. Παράλληλα, παρέχεται η δυνατότητα μαζικής ενημέρωσης (\en{broadcast}) από τον διοργανωτή προς όλους τους ενεργούς συμμετέχοντες.
    \item \textbf{Διαγραφή μηνυμάτων:} Εφαρμόζεται λογική διαγραφή μέσω των flags \en{deletedBySender} και \en{deletedByReceiver}. Το μήνυμα κρύβεται από τον χρήστη που το διέγραψε και η πραγματική διαγραφή από τη βάση δεδομένων εκτελείται αυτόματα μόνο όταν και οι δύο πλευρές επιλέξουν να το διαγράψουν.
\end{itemize}

\subsubsection{Γενικότερα}
\begin{itemize}[leftmargin=*]
    \item \textbf{Έλεγχος δικαιωμάτων:} Ενέργειες όπως τροποποίηση ή διαγραφή εκδηλώσεων προστατεύεται από μεθόδους ελέγχου όπως \en{isOrganizerOrAdmin} που επιτρέπει την πρόσβαση μόνο αν ο χρήστης είναι διαχειριστής ή ο διοργανωτής της εκδήλωσης. Αυτο επιβάλλεται και στην διαχείριση κρατήσεων και στην διαχείριση του προφιλ χρήστη.
    \item \textbf{Ασύγχρονα γεγονότα:} Για να μην καθυστερεί η απόκριση αιτήματος, σύνθετες ενέργειες όπως η ακύρωση κρατήσεων και η μαζική αποστολή ειδοποιήσεων λόγω ακύρωσης ενός event γίνονται ασύγχρονα. Το service δημοσιεύει ένα event μέσω του \en{ApplicationEventPublisher} και ένας listener αναλαμβάνει την εκτέλεση ασύγχρονα (σε άλλο thread) μετά την επιτυχή ολοκλήρωση της κύριας συναλλαγής/αιτήματος.
\end{itemize}







\subsection{Διαχείριση Πολυμέσων}
Για τη διαχείριση και την αποθήκευση των εικόνων χρησιμοποιείται η υπηρεσία \en{\textbf{FileSystemStorageService}}, η οποία αποθηκεύει ή διαγράφει αρχεία από τον προκαθορισμένο φάκελο αποθήκευσης στον δίσκο (default \en{media/ στον ριζικό κατάλογο της εφαρμογής}).

Όταν ανεβαίνει μια νέα εικόνα, η υπηρεσία δημιουργεί έναν ξεχωριστό υποφάκελο με βάση το \en{ID} της εκδήλωσης (\en{\texttt{media/\{eventId\}}}), «καθαρίζει» το όνομα του αρχείου για την αποφυγή επιθέσεων τύπου \en{path traversal} και το μετατρέπει σε μορφή \en{\textbf{.webp}}. Μέσω της βιβλιοθήκης \en{\textbf{Thumbnailator}}, η εικόνα περιορίζεται σε μέγιστη διάσταση \en{\textbf{1200x1200px}}, συμπιέζεται στο \en{\textbf{80\%}} της αρχικής της ποιότητας και αφαιρούνται όλα τα metadata για λόγους ασφαλείας.


\subsection{Διαχείριση Σφαλμάτων}
Όλα τα σφάλματα διαχειρίζονται μέσα από έναν \en{\textbf{GlobalExceptionHandler}}, ο οποίος μετατρέπει τα \en{domain exceptions} (\en{ValidationException, ResourceNotFoundException, AccessDeniedException} κ.ά.) σε τυποποιημένες \en{\textbf{RFC 7807} ProblemDetail JSON responses} με κατάλληλο \en{HTTP status} και επεξηγηματικό μήνυμα λάθους.

\newpage

% =========================================================
% ΚΕΦΑΛΑΙΟ 6: FRONTEND — REACT ΕΦΑΡΜΟΓΗ
% =========================================================
\section{Frontend — React Εφαρμογή}
Το \en{\textbf{Frontend}} είναι μια \en{\textbf{Single-Page Application} (SPA)} που υλοποιήθηκε με τη βιβλιοθήκη \en{\textbf{React}} και \en{\textbf{TypeScript}}. Η οργάνωση του κώδικα είναι λειτουργικά διαχωρισμένη σε φακέλους (\en{api, components, context, pages, styles} και \en{types}) για την καλύτερη διαχείριση και τη δομή του project.

\subsection{Routing και Δομή Σελίδων}
Η διαχείριση της πλοήγησης και των σελίδων γίνεται στην πλευρά του πελάτη (\en{client-side}) μέσω του \en{\textbf{React Router}} (\en{\texttt{react-router-dom}}). Όλο το περιεχομενό της εφαρμογής (\en{\texttt{AppContent}}) περιβάλλεται από τον \en{\texttt{AuthProvider}} για άμεση πρόσβαση στην κατάσταση σύνδεσης του χρήστη. Οι διαδρομές (\en{routes}) χωρίζονται σε τρεις βασικές κατηγορίες ανάλογα με τα δικαιώματα πρόσβασης:

\begin{itemize}[leftmargin=*, itemsep=6pt]
    \item \textbf{Δημόσιες Διαδρομές (\en{Public Routes}):} Είναι προσβάσιμες ελεύθερα από όλους τους χρήστες (επισκέπτες ή συνδεδεμένους).
        \begin{itemize}[leftmargin=0.6cm, itemsep=2pt]
            \item \en{\texttt{/}} : Η αρχική σελίδα καλωσορίσματος (\en{\texttt{WelcomePage}}) με τις προτεινόμενες εκδηλώσεις.
            \item \en{\texttt{/events}} : Σελίδα αναζήτησης και φιλτραρίσματος εκδηλώσεων (\en{\texttt{EventsPage}}).
            \item \en{\texttt{/events/:eventId}} : Αναλυτικές λεπτομέρειες εκδήλωσης με διαδραστικό χάρτη (\en{\texttt{EventDetailPage}}).
        \end{itemize}

    \item \textbf{Διαδρομές μόνο για Επισκέπτες (\en{Guest-Only Routes}):} Χρησιμοποιούν το φίλτρο \en{\texttt{GuestRoute}} και επιτρέπουν την πρόσβαση μόνο σε μη συνδεδεμένους χρήστες. Αν κάποιος έχει ήδη συνδεθεί, το σύστημα τον ανακατευθύνει αυτόματα στην αρχική σελίδα.
        \begin{itemize}[leftmargin=0.6cm, itemsep=2pt]
            \item \en{\texttt{/login}} : Σελίδα σύνδεσης στο σύστημα (\en{\texttt{LoginPage}}).
            \item \en{\texttt{/register}} : Σελίδα εγγραφής νέου λογαριασμού (\en{\texttt{RegisterPage}}).
            \item \en{\texttt{/pending-approval}} : Σελίδα ενημέρωσης για την εκκρεμή έγκριση του διαχειριστή (\en{\texttt{PendingApprovalPage}}).
        \end{itemize}

    \item \textbf{Προστατευμένες Διαδρομές (\en{Protected Routes}):} Χρησιμοποιούν το φίλτρο \en{\texttt{ProtectedRoute}} και απαιτούν ο χρήστης να έχει κάνει αυθεντικοποίηση με έγκυρο \en{JWT}, αλλιώς μεταφέρεται αυτόματα στη σελίδα σύνδεσης.
        \begin{itemize}[leftmargin=0.6cm, itemsep=2pt]
            \item \en{\texttt{/profile}} : Σελίδα προβολής και επεξεργασίας του προσωπικού προφίλ (\en{\texttt{ProfilePage}}).
            \item \en{\texttt{/admin}} : Ο κεντρικός πίνακας ελέγχου και διαχείρισης του διαχειριστή (\en{\texttt{AdminPage}}).
            \item \en{\texttt{/create-event}} : Φόρμα δημιουργίας νέας εκδήλωσης (\en{\texttt{CreateEventPage}}).
            \item \en{\texttt{/edit-event/:eventId}} : Σελίδα τροποποίησης στοιχείων εκδήλωσης από τον διοργανωτή μόνο (\en{\texttt{EditEventPage}}).
            \item \en{\texttt{/my-events}} : Πίνακας με τις εκδηλώσεις που διαχειρίζεται ο διοργανωτής (\en{\texttt{MyEventsPage}}).
            \item \en{\texttt{/my-events/:eventId/bookings}} : Λίστα με τις κρατήσεις χρηστών για μια συγκεκριμένη εκδήλωση από τον διοργανωτή μόνο (\en{\texttt{EventBookingsPage}}).
            \item \en{\texttt{/my-events/:eventId/broadcast}} : Σελίδα μαζικής αποστολής μηνυμάτων στους συμμετέχοντες από τον διοργανωτή μόνο (\en{\texttt{BroadcastMessagePage}}).
            \item \en{\texttt{/my-bookings}} : Πίνακας με τις κρατήσεις που έχει κάνει ο χρήστης (\en{\texttt{MyBookingsPage}}).
            \item \en{\texttt{/book/:eventId}} : Σελίδα επιλογής εισιτηρίων και ολοκλήρωσης κράτησης (\en{\texttt{BookingPage}}).
            \item \en{\texttt{/messages}} : Το \en{inbox} με τα μηνύματα και τις συνομιλίες (\en{\texttt{MessagePage}}).
            \item \en{\texttt{/messages/new}} : Σελίδα επιλογής παραλήπτη για έναρξη νέας συνομιλίας (\en{\texttt{NewMessagePage}}).
            \item \en{\texttt{/messages/compose}} : Φόρμα σύνταξης και αποστολής προσωπικού μηνύματος (\en{\texttt{ComposeMessagePage}}).
            \item \en{\texttt{/messages/:messageId}} : Προβολή των λεπτομερειών ενός συγκεκριμένου μηνύματος από τον αποστολέα ή παραλήπτη μόνο. (\en{\texttt{MessageDetailPage}}).
        \end{itemize}
\end{itemize}

Σε περίπτωση πρόσβασης κάποιας προστατευμένης διαδρομής από έναν συνδεδεμένο χρήστη που δεν έχει πρόσβαση (π.χ. σε \en{edit-event} διαδρομή με event στο οποίο δεν είναι ο διοργανωτής) τότε ο χρήστης μεταβαίνει σε προσαρμοσμένη σελίδα σφάλματος (\en{\texttt{ErrorPage}}) με κωδικό \textbf{\en{403 Unauthorized}}.

Παρόμοια, αν πληκτρολογηθεί μια διεύθυνση που δεν αντιστοιχεί σε καμία από τις παραπάνω ή αφορά \en{event/booking} που δεν υπάρχει τότε εμφανίζεται πάλι η σελίδα σφάλματος αλλά με κωδικό \textbf{\en{404 Not Found}}.

\subsection{Επιλεγμένες Σχεδιαστικές Αποφάσεις}
\begin{itemize}[label=\textbullet, leftmargin=*]
    \item \textbf{Αποκλειστικότητα Διαχειριστή (Admin Interceptor):} Αν συνδεθεί χρήστης με ρόλο διαχειριστή, το σύστημα τον ανακατευθύνει αμέσως στο \en{\texttt{/admin}} και του απαγορεύει την πλοήγηση στις σελίδες των απλών χρηστών.
    \item \textbf{Επιβεβαίωση Ενεργειών:} Στην ολοκλήρωση των κρατήσεων όπως και στην διαγραφή των εκδηλώσεων εμφανίζεται ένα επιβεβαιωτικό \en{dialog} πριν την οριστική αποστολή του αιτήματος, για την αποφυγή λανθασμένων ενεργειών.
    \item \textbf{Αυτόματη Συμπλήρωση Τοποθεσίας:} Χρήση του \en{\texttt{LocationAutocomplete}} component για την αυτόματη εύρεση και συμπλήρωση γεωγραφικών συντεταγμένων κατά τη δημιουργία μιας εκδήλωσης μέσω του \en{\textbf{Geoapify}}. Έχει απαγορευτεί η χειροκίνητη πληκτρολόγηση των διευθύνσεων, πόλης κ.τ.λπ. για την διευκολία στην εύρεση και αποθήκευση των γεογραφικών συντεταγμένων.
    \item \textbf{Κοινό Component Σελιδοποίησης:} Δημιουργήθηκε ένα ενιαίο \en{\texttt{Pagination}} component το οποίο επαναχρησιμοποιείται αυτούσιο σε όλες τις σελιδοποιημένες λίστες της εφαρμογής (εκδηλώσεις, κρατήσεις, εισερχόμενα).
    \item \textbf{Υποστήριξη Light/Dark Themes:} Πλήρης υποστήριξη εναλλαγής εμφάνισης με τη χρήση \en{CSS variables} για ομοιόμορφη εμπειρία χρήσης.
\end{itemize}

\newpage

% =========================================================
% ΚΕΦΑΛΑΙΟ 7: ΑΣΦΑΛΕΙΑ
% =========================================================
\section{Ασφάλεια}
Η ασφάλεια της πλατφόρμας υλοποιήθηκε σε όλα τα επίπεδα της εφαρμογής, εξασφαλίζοντας την ακεραιότητα των δεδομένων, την προστασία των χρηστών και την αποδοτική διαχείριση των πόρων του συστήματος.

\subsection{Κρυπτογραφημένη Επικοινωνία (SSL/TLS)}
Η κρυπτογραφημένη επικοινωνία εξασφαλίζεται στην πύλη εισόδου της εφαρμογής. Ο \en{Nginx Reverse Proxy} αναλαμβάνει εξολοκλήρου το \en{SSL termination} χρησιμοποιώντας ένα (\en{self-signed}) \en{SSL certificate}, προστατεύοντας όλη την εξωτερική κίνηση των χρηστών μέσω του πρωτοκόλλου \en{HTTPS}. Μετά την αποκρυπτογράφηση από τον \en{Nginx}, τα αιτήματα δρομολογούνται με ασφάλεια στο \en{Backend} εσωτερικά μέσω του απομονωμένου \en{Docker network} χρησιμοποιώντας το απλό πρωτόκολλο \en{HTTP}.

\subsection{Δικτυακή Απομόνωση (\en{Docker Networks})}
Για τη μέγιστη δυνατή ασφάλεια του συστήματος, τα \en{docker containers} επικοινωνούν μέσω δύο εντελώς ανεξάρτητων εσωτερικών δικτύων: το \en{\texttt{frontend-network}} και το \en{\texttt{backend-network}}. Η απομόνωση αυτή υλοποιείται ως εξής:
\begin{itemize}[leftmargin=*]
    \item \en{\textbf{Database}}: Ανήκει αποκλειστικά στο \en{\texttt{backend-network}} ώστε να επικοινωνεί μόνο με το \en{Backend}. Δεν εκθέτει καμία θύρα στο εξωτερικό σύστημα (\en{no open host ports}), αποκλείοντας οποιαδήποτε απευθείας πρόσβαση από το διαδίκτυο ή από το ίδιο το \en{Frontend}.
    \item \en{\textbf{Backend}}: Λειτουργεί ως γέφυρα μεταξύ των δύο δικτύων καθώς είναι συνδεδεμένο και στο \en{\texttt{frontend-network}} (για να δέχεται αιτήματα από τον \en{proxy}) και στο \en{\texttt{backend-network}} (για να μιλάει με τη βάση). Οι θύρες του προς τον έξω κόσμο είναι επίσης κλειστές ώστε να αποτρέπεται η απευθείας HTTP επικοινωνία μ'αυτό απο την θύρα του.
    \item \en{\textbf{Frontend}}: Είναι συνδεδεμένο στο \texttt{frontend-network} και είναι το μοναδικό τμήμα της εφαρμογής που εκθέτει θύρες στον εξωτερικό κόσμο (\en{\texttt{80}} και \en{\texttt{443}}). Ο \en{Nginx Reverse Proxy} δέχεται τα εξωτερικά αιτήματα και αναλαμβάνει να τα προωθήσει με ασφάλεια στο \en{backend} εσωτερικά.
\end{itemize}
Συνεπώς, κανένα εξωτερικό αίτημα δεν μπορεί να φτάσει στο \en{Backend} ή στη βάση δεδομένων χωρίς να περάσει υποχρεωτικά από τον έλεγχο και τη δρομολόγηση του \en{reverse proxy}. Ακόμη, και να υπάρξει security vulnerability στον Nginx που να επιτρέπει τον έλεγχο του με root δικαιώματα, δεν θα είναι, τουλάχιστον, εφικτή η πρόσβαση στην βάση δεδομένων.

\subsection{Αποθήκευση Κωδικών}
Για την προστασία των λογαριασμών, όλοι οι κωδικοί πρόσβασης των χρηστών αποθηκεύονται στη βάση δεδομένων αποκλειστικά σε κρυπτογραφημένη μορφή (\en{hashed}) χρησιμοποιώντας τον αλγόριθμο \en{BCrypt}. Το σύστημα έχει σχεδιαστεί έτσι ώστε οι κωδικοί να μην εκτίθενται ποτέ σε κανένα \en{DTO} ή \en{HTTP response}.

\subsection{JSON Web Tokens}
Η αυθεντικοποίηση στην εφαρμογή είναι πλήρως \en{stateless}. Κατά τη διαδικασία του \en{login}, εκδίδεται ένα \en{JWT (HS256)} με διάρκεια ζωής μίας ώρας και έπειτα, σε κάθε επόμενο αίτημα, το φίλτρο \en{\texttt{JwtAuthFilter}} ελέγχει την υπογραφή του \en{token} μέσω της βιβλιοθήκης \en{jjwt} και, εφόσον είναι έγκυρο, αναλαμβάνει αυτόματα τη φόρτωση των στοιχείων του χρήστη στο \en{SecurityContext} της εφαρμογής.

\subsection{Εξουσιοδότηση Βάσει Ρόλου}
Η διαχείριση των δικαιωμάτων γίνεται με σαφή διαχωρισμό τόσο στο \en{Security Config} όσο και σε επίπεδο \en{Controller}. Συγκεκριμένα, στο \en{\texttt{SecurityFilterChain}} ορίζονται οι καθολικοί κανόνες, δηλαδή ότι τα μονοπάτια \en{\texttt{/auth/**}} και οι αναζητήσεις εκδηλώσεων (\en{\texttt{GET /events/**}}) είναι ελεύθερα προσβάσιμα από όλους (\en{permitAll}). Για τις υπόλοιπες ενέργειες που απαιτείται αυθεντικοποίηση και έλεγχος των δικαιωμάτων, αυτό επιτυγχάνεται με τη χρήση του annotation \en{\texttt{@PreAuthorize}} στους \en{Controllers}.

\subsection{Προστασία από Spam / Σπατάλη Πόρων}
\begin{itemize}[leftmargin=*]
    \item \en{\textbf{CORS}:} Αυστηρά ρυθμισμένο στο \en{Spring Security} ώστε να δέχεται αιτήματα αποκλειστικά από το \en{URL} του \en{frontend}.
    \item \en{\textbf{XSS Protection}:} Η βιβλιοθήκη \en{React} προστατεύει την εφαρμογή από επιθέσεις \en{Cross-Site Scripting} κάνοντας αυτόματα \en{escape} οποιαδήποτε τιμή απεικονίζεται στο περιβάλλον εργασίας. Με αυτόν τον τρόπο, κακόβουλα κείμενα ή εντολές κώδικα που τυχόν προσπαθήσει να εισάγει ένας χρήστης αντιμετωπίζονται ως απλό κείμενο και δεν εκτελούνται ποτέ στον browser.
    \item \en{\textbf{Validation}:} Κάθε εισερχόμενο αίτημα ελέγχεται στην πλευρά του \en{Backend} με τη χρήση του \en{Jakarta Bean Validation} στα DTOs (π.χ. \en{\texttt{@NotBlank, @Min, @Email}}), απορρίπτοντας αμέσως άκυρα ή κακόβουλα δεδομένα πριν αυτά φτάσουν στη βάση.
\end{itemize}

\newpage

% =========================================================
% ΚΕΦΑΛΑΙΟ 8: ΕΞΑΓΩΓΗ ΕΚΔΗΛΩΣΕΩΝ ΣΕ XML / JSON
% =========================================================
\section{Εξαγωγή Εκδηλώσεων σε XML / JSON}
Ο διαχειριστής της εφαρμογής έχει τη δυνατότητα να εξάγει όλες τις καταχωρημένες εκδηλώσεις του συστήματος σε δομημένη μορφή αρχείου μέσω του \en{endpoint} \en{\texttt{/admin/export/events}} ορίζοντας την URL παράμετρο \en{\texttt{?format=xml}} ή \en{\texttt{?format=json}}. Η υλοποίηση αυτή βασίζεται σε ξεχωριστά, ειδικά σχεδιασμένα \en{Export DTOs} (\en{\texttt{EventsExport}}), τα οποία γίνονται \en{serialize} με τη χρήση των \en{\texttt{XmlMapper}} και \en{\texttt{JsonMapper}} της βιβλιοθήκης \en{Jackson}.

\subsection{Συμβατότητα με το DTD}
Για να εξασφαλιστεί η απόλυτη συμβατότητα με τους κανόνες και τη δομή του ζητούμενου \en{DTD} κατά την παραγωγή του \en{XML} αρχείου, χρησιμοποιούνται οι κατάλληλες \en{annotations} της \en{Jackson} πάνω από τα πεδία των \en{DTOs}:
\begin{itemize}[leftmargin=*]
    \item \en{\textbf{\texttt{@JsonPropertyOrder}}}: Χρησιμοποιείται για να επιβάλει μια αυστηρά καθορισμένη σειρά στα στοιχεία του \en{XML}, ώστε να μην αποτυγχάνει η επικύρωση λόγω ανακατέματος των πεδίων.
    \item \en{\textbf{\texttt{@JacksonXmlProperty(localName = "...")}}}: Μας επιτρέπει να αντιστοιχίσουμε τις ιδιότητες της \en{Java} με ονόματα στοιχείων που ξεκινούν με κεφαλαία γράμματα, ακολουθώντας τις ονοματολογίες του \en{DTD}.
    \item \en{\textbf{\texttt{@JacksonXmlProperty(isAttribute = true)}}}: Ορίζει ποια πεδία (όπως τα \en{\texttt{EventID, Latitude, Longitude, AvailableTickets}}) πρέπει να εμφανίζονται ως ιδιότητες (\en{attributes}) μέσα στο tag του στοιχείου και όχι ως εμφωλευμένα tags.
    \item \en{\textbf{\texttt{@JacksonXmlElementWrapper}}}: Αναλαμβάνει να δημιουργήσει τα απαραίτητα στοιχεία-φακέλους (\en{wrappers}) για τη σωστή οργάνωση των λιστών, όπως για παράδειγμα τα tags \en{\texttt{<TicketTypes>}} και \en{\texttt{<Bookings>}}.
    \item \en{\textbf{\texttt{@JsonInclude}}}: Ρυθμίζεται στην τιμή \en{\texttt{NON\_NULL}} ώστε αν κάποια πεδία (όπως οι φωτογραφίες μιας εκδήλωσης) είναι άδεια, να μην εμφανίζονται καθόλου στο τελικό αρχείο, αποτρέποντας τη δημιουργία άκυρων κενών tags.
\end{itemize}

Η \en{JSON} εκδοχή του αρχείου εξαγωγής παράγεται αυτόματα από τα ίδια ακριβώς \en{DTOs}, χρησιμοποιώντας τον έτοιμο \en{\texttt{JsonMapper}} που παρέχει το \en{Spring Boot} και τη μέθοδο \en{\texttt{writerWithDefaultPrettyPrinter()}} για βέλτιστη αναγνωσιμότητα.

\subsection{Λήψη Αρχείου}
Όποιο format και αν επιλεγεί, η απάντηση επιστρέφεται από το \en{Backend} με το κατάλληλο \en{Content-Type header} και με τη ρύθμιση \en{\texttt{Content-Disposition: attachment}} το οποίο αναγκάζει τον \en{browser} του χρήστη να ξεκινήσει την λήψη (\en{download}) του αρχείου στον υπολογιστή του διαχειριστή.

\newpage

% =========================================================
% ΚΕΦΑΛΑΙΟ 9: ΑΛΓΟΡΙΘΜΟΣ ΣΥΣΤΑΣΕΩΝ
% =========================================================

\section{Αλγόριθμος Συστάσεων}
Η εφαρμογή ενσωματώνει μια μηχανή προσωποποιημένων προτάσεων εκδηλώσεων για κάθε συνδεδεμένο χρήστη, βασισμένη στον αλγόριθμο \en{Biased Matrix Factorization (BMF)}. Ο αλγόριθμος αυτός αναλύει το ιστορικό αλληλεπιδράσεων που είναι αποθηκευμένο στον πίνακα \en{\texttt{User\_Event\_Interaction}} και προβλέπει το ενδιαφέρον ενός χρήστη για μελλοντικές εκδηλώσεις.

\subsection{Μαθηματικό Μοντέλο}
Η πρόβλεψη του σκορ προτίμησης $\hat{r}_{u,i}$ ενός χρήστη $u$ για μια εκδήλωση $i$ υπολογίζεται σύμφωνα με τη μαθηματική σχέση:
$$\hat{r}_{u,i} = \mu + b_u + b_i + \sum_{k=1}^{K} P_{u,k} \cdot Q_{i,k}$$

Όπου οι μεταβλητές ορίζονται ως εξής:
\begin{itemize}[leftmargin=*]
    \item $\mu$ : Η καθολική τάση του συστήματος (\en{global bias}), η οποία υπολογίζεται δυναμικά ως ο μέσος όρος όλων των αλληλεπιδράσεων.
    \item $b_u$ : Η προσωπική τάση του χρήστη (\en{user bias}), που αντικατοπτρίζει το πόσο θετικός ή αυστηρός είναι στις αλληλεπιδράσεις του.
    \item $b_i$ : Η τάση της εκδήλωσης (\en{item bias}), που δείχνει τη γενική δημοτικότητα της συγκεκριμένης εκδήλωσης στην πλατφόρμα.
    \item $P_u, Q_i$ : Τα διανύσματα των λανθανόντων παραγόντων (\en{latent factors}) για τον χρήστη και την εκδήλωση αντίστοιχα, τα οποία αποτυπώνουν τα κρυφά χαρακτηριστικά προτίμησης.
\end{itemize}

\subsection{Προετοιμασία Δεδομένων \& Offline Tuning (Jupyter Notebook)}
Πριν από την ενσωμάτωση του αλγορίθμου στο \en{Backend}, χρειάστηκε η εκτέλεση μιας offline φάσης προετοιμασίας σε \en{Jupyter Notebook} με το δοσμένο dataset. Στόχος της διαδικασίας αυτής ήταν ο καθαρισμός/προετοιμασία του αρχικού \en{dataset} και η εύρεση των ιδανικών υπερπαραμέτρων εκπαίδευσης για την μέγιστη ακρίβεια στις συστάσεις. Η offline αυτή διαδικασία βασίστηκε στα εξής βήματα:
\begin{itemize}[leftmargin=*]
    \item \textbf{Κατασκευή και εμπλουτισμός Ratings:} Φορτώθηκαν τα αρχικά δεδομένα από τα αρχεία \en{event\_interest.csv} και \en{event\_attendees.csv}. Τα σήματα ενδιαφέροντος μετατράπηκαν σε βαθμολογίες (\en{Interested} = 1.0, \en{Not Interested} = 0.0) και εμπλουτίστηκαν με τα δεδομένα προσέλευσης (\en{Attendance}), όπου το \en{yes} αντιστοιχεί σε 1.0, το \en{maybe} σε 0.6 και το \en{no} σε 0.0.
    \item \textbf{Καθαρισμός δεδομένων:} Για την εξάλειψη του θορύβου, φιλτραρίστηκαν και αφαιρέθηκαν χρήστες και εκδηλώσεις με λιγότερες από 3 αλληλεπιδράσεις. Στη συνέχεια, έγινε \en{remap} των \en{IDs} των χρηστών και των εκδηλώσεων σε μια συνεχή λίστα ακεραίων (0 έως n-1) για τη βέλτιστη διαχείριση των πινάκων.
    \item \textbf{Αξιολόγηση και \en{Grid Search}:} Εφαρμόστηκε η μέθοδος \en{5-Fold Cross-Validation} για την αντικειμενική αξιολόγηση του μοντέλου \en{from scratch}. Μέσω \en{Grid Search} δοκιμάστηκαν διαφορετικοί συνδυασμοί για τον αριθμό των παραγόντων, το ρυθμό μάθησης και την κανονικοποίηση, μετρώντας το σφάλμα \en{RMSE}, καθώς και τις μετρικές κατάταξης \en{Precision@10} και \en{Recall@10}.
\end{itemize}
Η διαδικασία αυτή έβγαλε ιδανικό \en{config} το εξής: 3 λανθάνοντες παράγοντες, 0.025 learning rate, 0.03 regularization και 60 epoch. Οι ρυθμίσεις αυτές μεταφέρθηκαν αυτούσιες στον πίνακα \en{\texttt{Recommendation\_Config}} της βάσης δεδομένων για τα μετέπειτα \en{trainings} με αληθινούς χρήστες και εκδηλώσεις στην εφαρμογή.

\subsection{Καταγραφή Δεδομένων \& Περιορισμοί στο Backend}
Η συλλογή των νέων δεδομένων αλληλεπίδρασης γίνεται πλέον σε πραγματικό χρόνο από το \en{Backend} ανάλογα με τη δραστηριότητα των συνδεδεμένων χρηστών, δίνοντας σκορ 0.4 για την απλή προβολή μιας εκδήλωσης και 1.0 για την ολοκλήρωση μιας κράτησης.

Για την αποθήκευση αυτών των δεδομένων, το \en{interaction repository} εκτελεί ένα \en{SQL upsert} ερώτημα (\en{\texttt{INSERT ... ON DUPLICATE KEY UPDATE}}), το οποίο επιβάλλει δύο κρίσιμους κανόνες:
\begin{itemize}[leftmargin=*]
    \item \textbf{Προτεραιότητα ενέργειας:} Χρησιμοποιεί τη συνάρτηση \en{\texttt{GREATEST(rating, :rating)}}, διασφαλίζοντας ότι αν ένας χρήστης έχει ήδη δει μια εκδήλωση (0.4) και μετά κάνει κράτηση (1.0), το σκορ θα αναβαθμιστεί, ενώ αν συμβεί το αντίθετο, δεν θα υποβαθμιστεί η αξία της κράτησης.
    \item \textbf{Μπλοκάρισμα διοργανωτή:} Μέσω του φίλτρου \en{\texttt{WHERE e.organizer\_id != :userId}}, απαγορεύεται ρητά η καταγραφή αλληλεπιδράσεων όταν ένας χρήστης αλληλεπιδρά με δικές του εκδηλώσεις.
\end{itemize}

\subsection{Αυτοματοποιημένο Pipeline Εκπαίδευσης}
Η επανεκπαίδευση του μοντέλου είναι πλήρως αυτοματοποιημένη και εκτελείται στο παρασκήνιο κάθε 30 λεπτά μέσω ενός background scheduler (\en{\texttt{@Scheduled}}). Η διαδικασία εκτελείται ασύγχρονα (\en{\texttt{@Async}}) σε ξεχωριστό \en{thread pool} προκειμένου να μην επιβαρύνει καθόλου τις \en{HTTP} απαντήσεις των χρηστών. Με την έναρξη, το \en{service} εκκαθαρίζει άμεσα την \en{cache} των προτεινόμενων εκδηλώσεων για τους επισκέπτες, υπολογίζει δυναμικά το νέο \en{global bias} ως τον τρέχοντα μέσο όρο όλων των εγγραφών της βάσης, και ξεκινά τη διαδικασία βελτιστοποίησης με τον αλγόριθμο \en{Stochastic Gradient Descent (SGD)}.

\subsection{Βελτιστοποιήσεις Μνήμης \& Εκτέλεση SGD}
Για την επίτευξη υψηλών επιδόσεων κατά τη διάρκεια των 60 εποχών εκπαίδευσης, χρησιμοποιούνται διάφορες τεχνικές:
\begin{itemize}[leftmargin=*]
    \item \textbf{Caching σε HashMaps:} Πριν ξεκινήσουν οι επαναλήψεις, όλες οι υπάρχουσες εγγραφές των χρηστών και των εκδηλώσεων φορτώνονται από τη βάση σε τοπικά \en{HashMaps}. Αυτό επιτρέπει ακαριαίες αναζητήσεις (\en{O(1)} lookups) και ενημερώσεις των διανυσμάτων στη μνήμη, αποφεύγοντας χιλιάδες ενδιάμεσα queries.
    \item \textbf{Αποφυγή Ordering Bias:} Στην αρχή κάθε εποχής, εκτελείται η μέθοδος \en{\texttt{Collections.shuffle(interactions)}}. Η τυχαία αναδιάταξη των αλληλεπιδράσεων είναι απαραίτητη για την ορθή λειτουργία του \en{SGD}, καθώς αποτρέπει το μοντέλο από το να εγκλωβιστεί σε τοπικά ελάχιστα λόγω της σειράς εισαγωγής των δεδομένων.
    \item \textbf{Δυναμική αρχικοποίηση:} Αν εντοπιστεί νέος χρήστης ή νέα εκδήλωση τότε δημιουργείται το διάνυσμα του στη μνήμη σε πραγματικό χρόνο. Οι λανθάνοντες παράγοντες αρχικοποιούνται με τυχαίες μικρές τιμές και το \en{bias} στο 0.0.
\end{itemize}
Σε κάθε βήμα, ο αλγόριθμος συγκρίνει το πραγματικό σκορ με την πρόβλεψη, υπολογίζει το σφάλμα (\en{error}) και εφαρμόζει τους κανόνες ενημέρωσης του \en{SGD} τροποποιώντας τις τιμές στα τοπικά αντικείμενα. Μόλις ολοκληρωθούν όλες οι εποχές, τα ενημερωμένα διανυσματα αποθηκεύονται μαζικά πίσω στη βάση δεδομένων.


\subsection{Real-time Λειτουργία Συστάσεων}
Για την εμφάνιση των προσωποποιημένων προτάσεων και των αποτελεσμάτων αναζήτησης, οι υπολογισμοί του μοντέλου εκτελούνται σε πραγματικό χρόνο απευθείας στο επίπεδο της βάσης δεδομένων. Στο \en{EventRepository}, η ανάκτηση των εκδηλώσεων υλοποιείται με custom JPQL queries που πραγματοποιούν \en{LEFT JOIN} με τους πίνακες των vectors (\en{UserRecommendationVector} και \en{EventRecommendationVector}), εφαρμόζοντας μια αυστηρή ιεραρχία ταξινόμησης (\en{ORDER BY}):

\begin{enumerate}[leftmargin=*]
    \item \textbf{Κατάσταση Εκδήλωσης:} Οι εκδηλώσεις ομαδοποιούνται με βάση το status τους, δίνοντας απόλυτη προτεραιότητα στις ενεργές (\en{PUBLISHED}), ακολουθούμενες από τις ολοκληρωμένες (\en{COMPLETED}) και τέλος τις ακυρωμένες (\en{CANCELLED}).
    \item \textbf{Σκόρ Σύστασης:} Μέσα σε κάθε ομάδα κατάστασης, η κατάταξη γίνεται φθίνουσα με βάση το μαθηματικό γινόμενο των λανθανόντων παραγόντων και του item bias. Σημειώνεται ότι ο καθολικός μέσος όρος ($\mu$) παραλείπεται από το query καθώς αποτελεί σταθερά και δεν επηρεάζει τη σχετική διάταξη των αποτελεσμάτων.
    \item \textbf{Χρονος Έναρξης:} Σε περίπτωση ισοβαθμίας στα σκορ συστάσεων, οι εκδηλώσεις ταξινομούνται με βάση την ημερομηνία έναρξης (\en{startDatetime}) σε αύξουσα σειρά, εμφανίζοντας πρώτα τις πιο κοντινές χρονικά.
\end{enumerate}

Παράλληλα, το query αντιμετωπίζει δυναμικά τις περιπτώσεις ελλιπών δεδομένων μέσω της συνάρτησης \en{COALESCE}, η οποία λειτουργεί ως δικλείδα ασφαλείας (\en{fallback}) σε δύο κρίσιμα σενάρια:
\begin{itemize}[leftmargin=*]
    \item \textbf{Επισκέπτες:} Όταν το \en{userId} είναι \en{null}, το \en{LEFT JOIN} στον πίνακα του χρήστη επιστρέφει κενές τιμές. Η \en{COALESCE} μετατρέπει τα \en{nulls} των παραγόντων του χρήστη σε 0.0, εκμηδενίζοντας το εσωτερικό γινόμενο. Έτσι, ο τύπος απλοποιείται αυτόματα στο \en{coalesce(ev.bias, 0)}, επιστρέφοντας συστάσεις βασισμένες αποκλειστικά στη γενική δημοτικότητα (bias) της κάθε εκδήλωσης.
    \item \textbf{Νέες Οντότητες (\en{Cold Start}):} Για ολοκαίνουργιους χρήστες ή νέα events που δεν έχουν προλάβει να συμπεριληφθούν σε κάποια εκπαίδευση του scheduler, η \en{COALESCE} αποτρέπει την κατάρρευση του query επιστρέφοντας μια ουδέτερη τιμή (0.0), κατατάσσοντας τις οντότητες αυτές σε μια ουδέτερη βάση χωρίς να αλλοιώνει την υπόλοιπη στοίχιση.
\end{itemize}


\newpage

% =========================================================
% ΚΕΦΑΛΑΙΟ 10: ΟΔΗΓΙΕΣ ΕΓΚΑΤΑΣΤΑΣΗΣ & ΕΚΤΕΛΕΣΗΣ
% =========================================================
\section{Οδηγίες Εγκατάστασης \& Εκτέλεσης}

\subsection{Προαπαιτούμενα}
\begin{itemize}
    \item \textbf{\href{https://www.docker.com/products/docker-desktop/}{Docker Desktop}} (installed and running)
\end{itemize}

\subsection{Τοπική Εγκατάσταση \& Εκτέλεση}

\begin{enumerate}
    \item Ανοίξτε ένα τερματικό στον ριζικό κατάλογο της εφαρμογής (\en{\texttt{Planent\_Web\_App/}})
    \item Εκτελέστε την ακόλουθη εντολή για να ξεκινήσουν όλες τις υπηρεσίες (\en{docker containers}):
\end{enumerate}

\begin{lstlisting}[language=bash]
docker compose up --build
\end{lstlisting}
Αυτό κατασκευάζει και εκκινεί αυτόματα τα containers της \en{\textbf{Βάσης Δεδομένων}}, του \en{\textbf{Backend}} και του \en{\textbf{Frontend}} μέσα σε απομονωμένα εσωτερικά δίκτυα \en{Docker}.

\subsection{Πρόσβαση στην Εφαρμογή}

Μόλις όλα τα containers τεθούν σε λειτουργία, η εφαρμογή (ιστότοπος) θα είναι διαθέσιμη στη διεύθυνση \textbf{\href{https://localhost}{https://localhost}}.

Η εφαρμογή επιβάλλει τη χρήση HTTPS χρησιμοποιώντας ένα τοπικό αυτο-υπογεγραμμένο (self-signed) πιστοποιητικό SSL, επομένως ο περιηγητής (browser) θα εμφανίσει μια προειδοποίηση "Η σύνδεσή σας δεν είναι ιδιωτική" (``Your connection is not private''). Απλώς κάντε κλικ στο \texttt{Advanced} (Για προχωρημένους) και στη συνέχεια στο \texttt{Proceed to localhost (unsafe)} (Συνέχεια σε localhost).

Ένας προεπιλεγμένος χρήστης διαχειριστή (admin) δημιουργείται επίσης αυτόματα με όνομα χρήστη και κωδικό πρόσβασης \texttt{admin} (μπορεί να αλλάξει στο \texttt{backend/src/main/java/com/uoa/planent/PlanentApplication\#initDatabase}).


\subsection{Τερματισμός \& Επαναφορά:}
\begin{lstlisting}[language=bash]
docker compose down
\end{lstlisting}
Όλα τα αποθηκευμένα δεδομένα διατηρούνται με ασφάλεια στον volume \texttt{mysql-data} που δημιουργήθηκε από το container της \textbf{Βάσης Δεδομένων}, ακόμη και μετά τον τερματισμό.
Εκτελέστε ξανά την εφαρμογή με \texttt{docker compose up}, ή ξανά με flag \texttt{--build} εάν έχουν αλλάξει αρχεία.


\subsection{Καθαρισμός \& Διαγραφή:}
Για να αφαιρέσετε τα docker αρχεία που δημιουργήθηκαν για αυτήν την εφαρμογή, χρησιμοποιήστε μία από τις ακόλουθες εντολές:

\vspace{1\baselineskip}
\noindent \\ \\ \textbf{1. Μόνο Containers \& Images (χωρίς το volume της Βάσης Δεδομένων):}
\begin{lstlisting}[language=bash]
docker compose down --rmi all
\end{lstlisting}

\vspace{1\baselineskip}
\noindent \textbf{2. Τα πάντα (μαζι με το volume της Βάσης Δεδομένων):}
\begin{lstlisting}[language=bash]
docker compose down -v --rmi all
\end{lstlisting}


\newpage

% =========================================================
% ΚΕΦΑΛΑΙΟ 11: ΕΠΙΛΟΓΟΣ
% =========================================================
\section{Επίλογος}
Η ανάπτυξη της εφαρμογής \en{Planent} αποτέλεσε μια ολοκληρωμένη άσκηση σχεδίασης και υλοποίησης σύγχρονης διαδικτυακής εφαρμογής, με χρήση επαγγελματικών τεχνολογιών.

Κατά τη διάρκεια της υλοποίησης αντιμετωπίστηκαν αρκετές προκλήσεις:
\begin{itemize}[label=\textbullet]

\item \textbf{Κοινό περιβάλλον συνεργασίας:} Μία από τις πρώτες προκλήσεις ήταν να μπορούμε να εργαζόμαστε και οι δύο ταυτόχρονα πάνω στα ίδια δεδομένα. Για τον σκοπό αυτό, κατά τη φάση της ανάπτυξης, στήσαμε και ενοικιάσαμε μια κοινή (\en{hosted}) \en{MySQL} βάση δεδομένων στην πλατφόρμα \en{Railway}, ώστε να έχουμε κοινή πρόσβαση και να συγχρονίζουμε εύκολα τη δουλειά μας. Παράλληλα, χρησιμοποιήσαμε σύστημα ελέγχου εκδόσεων (\en{version control}) μέσω της πλατφόρμας \en{GitHub}, όπου δημιουργήσαμε το αντίστοιχο \en{repository} και φροντίζαμε να συγχρονίζουμε τακτικά τις ενημερώσεις (\en{updates}) μας, ώστε ο καθένας να συμπληρώνει τη δουλειά του άλλου.

\item \textbf{Επιλογή τεχνολογικής στοίβας \& setup:} Σημαντικός προβληματισμός υπήρξε η τελική επιλογή των τεχνολογιών και του τρόπου στησίματος του \en{Backend}. Αρχικά ξεκινήσαμε την υλοποίηση σε \en{TypeScript}, ωστόσο στην πορεία αποφασίσαμε να μεταβούμε σε \en{Spring Boot} με \en{Java}, καθώς μας προσέφερε μεγαλύτερη δομή, ωριμότητα και ευκολία στη διαχείριση της ασφάλειας και της επιμονής των δεδομένων (\en{persistence}).

\item \textbf{Σχεδίαση του σχήματος της βάσης:} Η μοντελοποίηση όλων των απαραίτητων πινάκων και των μεταξύ τους σχέσεων αποτέλεσε μια απαιτητική διαδικασία, ώστε να καλύπτονται πλήρως οι λειτουργικές απαιτήσεις, οι περιορισμοί ακεραιότητας, αλλά και η συμβατότητα με το ζητούμενο \en{DTD}.

\item \textbf{Ασύγχρονες λειτουργίες:} Η εκτέλεση χρονοβόρων ενεργειών (όπως η μαζική αποστολή μηνυμάτων κατά την ακύρωση μιας εκδήλωσης) χωρίς να καθυστερεί η απόκριση του \en{API} απαίτησε τη χρήση \en{Spring Application Events} και \en{@TransactionalEventListener} στη φάση \en{AFTER\_COMMIT}.

\item \textbf{Συμμόρφωση εξαγωγής με το \en{DTD}:} Η παραγωγή ενός \en{XML} αρχείου που να ταυτίζεται ακριβώς με τη δομή και τη σειρά των στοιχείων του ζητούμενου \en{DTD} απαίτησε προσεκτική χρήση των \en{annotations} της \en{Jackson} (\en{\texttt{@JsonPropertyOrder}}, \en{\texttt{@JacksonXmlProperty}} κ.ά.).

\item \textbf{Απόδοση \& Εμπειρία Χρήστη (\en{UX}):} Δόθηκε ιδιαίτερη προσπάθεια ώστε οι αποκρίσεις του συστήματος να είναι όσο το δυνατόν πιο γρήγορες και η εμπειρία χρήσης πιο καθαρή και ομαλή. Αυτό απαίτησε βελτιστοποιήσεις τόσο στο \en{Backend} (π.χ. \en{batch fetching} για την αποφυγή του προβλήματος \en{N+1}, χρήση \en{cache}) όσο και στο \en{Frontend} (καθαρό \en{UI}, σελιδοποίηση, ασύγχρονη φόρτωση δεδομένων).

\item \textbf{Εκπαίδευση \& ενσωμάτωση του αλγορίθμου συστάσεων:} Από τα πιο σύνθετα κομμάτια ήταν η υλοποίηση του \en{Biased Matrix Factorization} από το μηδέν, η εύρεση των κατάλληλων υπερπαραμέτρων μέσω \en{offline tuning}, η ενσωμάτωσή του στην εφαρμογή και, τέλος, οι δοκιμές με πραγματικές εκδηλώσεις για την επαλήθευση της ορθής λειτουργίας των προτάσεων.
\end{itemize}

Ο αλγόριθμος συστάσεων ενσωματώθηκε πλήρως στην εφαρμογή, προσφέροντας μια προσωποποιημένη εμπειρία χρήσης βασισμένη στις πραγματικές ανάγκες και προτιμήσεις των επισκεπτών της πλατφόρμας.

\end{document}